% =============================================================================
%  bridge.tex — From the Bloch Sphere to Cosmological Posteriors
%  A code-level bridge between classical and quantum Bayesian inference.
%
%  Audience: a reader fluent in cosmology and Bayesian statistics whose
%  quantum-computing background does not exceed the Bloch sphere. Every
%  quantum construction is introduced FROM a single qubit and TOWARD its
%  algorithmic role, and every code function is anchored to one of three
%  categories (Physics / Statistics / Quantum Computation) with its
%  time/space complexity.
%
%  Compile:  pdflatex bridge.tex  (run twice for the table of contents)
% =============================================================================
\documentclass[11pt]{article}

\usepackage[utf8]{inputenc}
\usepackage[T1]{fontenc}
\usepackage[margin=1in]{geometry}
\usepackage{amsmath,amssymb,amsthm}
\usepackage{mathtools}
\usepackage{braket}            % Dirac notation: \ket, \bra, \braket
\usepackage{booktabs}
\usepackage{longtable}
\usepackage{array}
\usepackage{xcolor}
\usepackage{listings}
\usepackage{seqsplit}
\usepackage{hyperref}
\hypersetup{colorlinks=true, linkcolor=blue!55!black,
            citecolor=blue!55!black, urlcolor=blue!55!black}
\usepackage{enumitem}
\usepackage{tikz}
\usetikzlibrary{quantikz2} % quantum circuits; falls back gracefully if absent

% --- code listing style ------------------------------------------------------
\definecolor{codebg}{gray}{0.96}
\definecolor{codekw}{rgb}{0.0,0.3,0.6}
\definecolor{codecm}{rgb}{0.3,0.5,0.3}
\lstdefinestyle{py}{
  language=Python, backgroundcolor=\color{codebg},
  basicstyle=\ttfamily\footnotesize, keywordstyle=\color{codekw}\bfseries,
  commentstyle=\color{codecm}\itshape, stringstyle=\color{red!50!black},
  showstringspaces=false, breaklines=true, frame=single,
  framerule=0pt, columns=fullflexible, xleftmargin=4pt, xrightmargin=4pt}
\lstset{style=py}

% --- theorem-like environments ----------------------------------------------
\newtheorem{definition}{Definition}[section]
\newtheorem{proposition}{Proposition}[section]
\newtheorem{remark}{Remark}[section]

% --- shorthands --------------------------------------------------------------
\newcommand{\Om}{\Omega_m}
\newcommand{\OL}{\Omega_\Lambda}
\newcommand{\Or}{\Omega_r}
\newcommand{\ODE}{\Omega_{\mathrm{DE}}}
\newcommand{\Ez}{E(z;\theta)}
\newcommand{\Ezsq}{E^2(z;\theta)}
\newcommand{\code}[1]{\texttt{#1}}
\newcommand{\bigO}{\mathcal{O}}
\newcommand{\KL}{\mathrm{KL}}
\newcommand{\Rhat}{\hat{R}}
\newcommand{\faithful}{\textsc{faithful}}
\newcommand{\algorithmic}{\textsc{algorithmic}}

\title{\bfseries From the Bloch Sphere to Cosmological Posteriors\\[2pt]
\large A Code-Level Bridge Between Classical and Quantum Bayesian Inference}
\author{Andr\'es Garc\'ia Rivera\\
\small Engineering Physics, Universidad Iberoamericana (IBERO)}
\date{\today}

\begin{document}
\maketitle

\begin{abstract}
\noindent
This document is the theoretical companion to the \emph{Quantum Algorithms
for Cosmology} codebase. It is written for a reader who is at home with
cosmology and Bayesian statistics but whose quantum-computing knowledge stops
at the Bloch sphere. Its purpose is threefold: to separate and explain every
function of the code under three headings (Physics, Statistics, Quantum
Computation); to build a bridge from a single qubit to the quantum samplers,
so the connection between circuit and inference is explicit; and to give the
time and space complexity of each component. The guiding scientific principle
throughout is \emph{fidelity, not advantage}: for classical parametric model
fitting there is no structural quantum speed-up to be had, so the genuine
contribution is to show that quantum methods faithfully \emph{replicate}
trusted classical results. We make that idea precise by casting the whole
project as a \emph{classical--quantum ablation framework}, in which each
algorithmic component is a substitution point that is either \faithful{}
(a quantum re-encoding of the exact classical rule, expected to reproduce it)
or \algorithmic{} (a genuinely different quantum algorithm, expected to
differ).
\end{abstract}

\tableofcontents
\newpage

% =============================================================================
\section{Preliminaries}
% =============================================================================

\subsection{Purpose and the code\texorpdfstring{$\to$}{->}theory map}

The codebase is organized so that \emph{physics}, \emph{statistics}, and
\emph{sampling/optimization} never bleed into one another. All physics lives
in \code{cosmo\_core.py}; the samplers in \code{cosmo\_modular\_quantum.py}
(simulator) and \code{qpu\_cosmo\_samplers.py} (hardware); the global
optimizers in \code{cosmo\_genetic\_optimizers.py}. This document mirrors that
separation. Table~\ref{tab:map} is the master index: each function appears
under exactly one category, with the section that explains it and its
complexity.

\begin{table}[h]
\centering
\footnotesize
\begin{tabular}{@{}lllc@{}}
\toprule
\textbf{Function / object} & \textbf{Category} & \textbf{Section} & \textbf{Complexity (time)}\\
\midrule
\code{CosmoModel}, \code{\_E2\_*}        & Physics    & \S\ref{sec:physics}  & $\bigO(N_z)$ \\
\code{Posterior.\_mu\_theory}            & Physics    & \S\ref{sec:physics}  & $\bigO(N_z)$ \\
\code{Posterior.log\_prob\_batch}        & Statistics & \S\ref{sec:stats}    & $\bigO(B(N_z{+}N_{\mathrm{SN}}^2))$ \\
\code{chi2\_pantheon\_plus}              & Statistics & \S\ref{sec:stats}    & $\bigO(N_{\mathrm{SN}}^2)$ \\
\code{autocorr\_time\_fft}               & Statistics & \S\ref{sec:stats}    & $\bigO(N\log N)$ \\
\code{split\_rhat}                       & Statistics & \S\ref{sec:stats}    & $\bigO(MN)$ \\
\code{QMCMCModular}                      & Statistics/QC & \S\ref{sec:qmcmc} & per step $\bigO(M\cdot c_{\mathrm{post}})$ \\
\code{hadamard\_accept\_log\_batch}      & QC (\faithful) & \S\ref{sec:bridge} & $\bigO(M)$ \\
\code{QuantumProposalEngine}             & QC (\algorithmic) & \S\ref{sec:bridge} & $\bigO(2^{n_q})$/refill \\
\code{QVMCModular}                       & Statistics/QC & \S\ref{sec:qvmc}  & target $\bigO(2^{n\cdot\mathrm{nqpp}})$ \\
\code{estimate\_grid\_window}            & Statistics & \S\ref{sec:qvmc}     & $\bigO(n_{\mathrm{steps}}M c_{\mathrm{post}})$ \\
\code{CGA}/\code{QGA}                    & Statistics/QC & \S\ref{sec:genetic} & $\bigO(G\cdot P\cdot d)$ \\
\code{compute\_quantumness}              & Framework  & \S\ref{sec:ablation} & $\bigO(1)$ \\
\bottomrule
\end{tabular}
\caption{Master map: every component, its category, the section that explains
it, and its dominant time complexity. $B$ = batch size, $N_z$ = distance-grid
size, $N_{\mathrm{SN}}$ = number of supernovae, $M$ = chains, $N$ = chain
length, $d$ = parameters, $\mathrm{nqpp}$ = qubits per parameter, $G$ =
generations, $P$ = population, $c_{\mathrm{post}}$ = cost of one posterior
evaluation.}
\label{tab:map}
\end{table}

\subsection{Notation and the guiding principle}

We write the dimensionless Hubble function as $\Ez$ with $\Ezsq =
H^2(z;\theta)/H_0^2$, and $\theta$ for the parameter vector (by convention
$\theta_0=\Om$, $\theta_1=H_0$). Quantum states use Dirac notation; a single
qubit is $\ket{\psi}=\alpha\ket{0}+\beta\ket{1}$ with $|\alpha|^2+|\beta|^2=1$.

\begin{remark}[Fidelity, not advantage]
For fitting a handful of cosmological parameters to fixed data, the bottleneck
is the \emph{likelihood evaluation}, which is classical arithmetic over the
data. No quantum subroutine in this project changes that asymptotics, and we
make no claim of quantum advantage. The scientific value is the demonstration
that quantum re-encodings of the inference steps \emph{reproduce} the
classical posterior --- a fidelity result. Section~\ref{sec:ablation} turns
this into a falsifiable framework.
\end{remark}

% =============================================================================
\section{Physics}
\label{sec:physics}
% =============================================================================

The only physics a model must expose is $\Ezsq$; every observable is derived
generically from it. This is the \code{CosmoModel} contract.

\subsection{The Friedmann function and the model registry}

\begin{definition}[Background model]
A flat FLRW background with matter, radiation and dark energy obeys
\begin{equation}
\Ezsq = \Om(1+z)^3 + \Or(1+z)^4 + \ODE\, f_{\mathrm{DE}}(z),
\end{equation}
with $\ODE = 1-\Om-\Or$ (flatness) and $f_{\mathrm{DE}}(z)$ the dark-energy
density evolution, $f_{\mathrm{DE}}(0)=1$. The Hubble rate is
$H(z;\theta)=H_0\,\Ez$.
\end{definition}

The five registered models differ only in $f_{\mathrm{DE}}$:
\begin{align}
\text{$\Lambda$CDM:}\quad & f_{\mathrm{DE}}=1,\\
\text{wCDM:}\quad & f_{\mathrm{DE}}=(1+z)^{3(1+w)},\\
\text{CPL:}\quad & f_{\mathrm{DE}}=(1+z)^{3(1+w_0+w_a)}\exp\!\big(\!-3w_a z/(1+z)\big),\\
\text{PEDE:}\quad & f_{\mathrm{DE}}=1-\tanh\!\big(\log_{10}(1+z)\big),\\
\text{GEDE:}\quad & f_{\mathrm{DE}}=\frac{1-\tanh\!\big(\Delta\log_{10}\frac{1+z}{1+z_t}\big)}
{1+\tanh\!\big(\Delta\log_{10}(1+z_t)\big)},\quad (1+z_t)=\Big(\tfrac{1-\Om}{\Om}\Big)^{1/3}.
\end{align}

\begin{proposition}[Emergent dark-energy normalization]
PEDE and GEDE satisfy $f_{\mathrm{DE}}(0)=1$ exactly, hence $\Ezsq|_{z=0}=1$.
\end{proposition}
\begin{proof}
For PEDE, $\tanh(\log_{10}1)=\tanh 0=0$, so $f_{\mathrm{DE}}(0)=1$. For GEDE at
$z=0$, $\tfrac{1+z}{1+z_t}=\tfrac{1}{1+z_t}$, so the numerator is
$1-\tanh(-\Delta\log_{10}(1+z_t))=1+\tanh(\Delta\log_{10}(1+z_t))$, identical
to the denominator; their ratio is $1$. This holds for every $\Om,\Delta$ and
is asserted in \code{\seqsplit{tests/test\_core.py::test\_emergent\_de\_normalization}}.
\end{proof}

\begin{remark}[Convention caveat checked at $z\neq0$]
An inverted sign of $z_t$ would also give $f_{\mathrm{DE}}(0)=1$ but the
opposite high-$z$ evolution. The test
\code{\seqsplit{test\_gede\_reduces\_to\_lcdm\_at\_delta\_zero}} therefore checks
$\Delta=0\Rightarrow f_{\mathrm{DE}}\equiv1$ at $z=0.5,1,2$, not only at $z=0$.
\end{remark}

\subsection{From $\Ezsq$ to observables}

Two observables enter the likelihood. The Hubble rate $H(z;\theta)$ is read
directly. The distance modulus requires the comoving distance:
\begin{equation}
d_C(z) = \frac{c}{H_0}\int_0^z \frac{dz'}{\Ez[z']},\qquad
d_L=(1+z)\,d_C,\qquad
\mu_{\mathrm{th}}=5\log_{10}\!\big(d_L/\text{Mpc}\big)+25.
\end{equation}

\begin{remark}[Flatness, made explicit]
The code takes the transverse comoving distance equal to $d_C$, valid only for
$\Omega_k=0$. All five models are flat. A curved extension would need
$d_M=(c/H_0\sqrt{|\Omega_k|})\,\mathrm{sinn}(\sqrt{|\Omega_k|}H_0 d_C/c)$;
\code{\_mu\_theory} documents this so a future non-flat model is not silently
wrong.
\end{remark}

\paragraph{Complexity.}
The integral is computed once on a fixed grid of $N_z$ points by cumulative
trapezoid, then interpolated to the supernova redshifts: $\bigO(N_z)$ time and
space, and crucially \emph{model-agnostic} --- the same code serves CPL or
GEDE. The earlier point-by-point \code{scipy.integrate.quad} was $\bigO(N_{
\mathrm{SN}}\cdot\text{(integrand evals)})$ and $\sim$10$\times$ slower.

% =============================================================================
\section{Statistics}
\label{sec:stats}
% =============================================================================

\subsection{Posterior, likelihoods and the marginalized $M_{\mathrm{abs}}$}

The log-posterior ties model, data and prior:
\begin{equation}
\log P(\theta\mid \mathrm{data}) = \log\pi(\theta) - \tfrac12\chi^2(\theta),
\qquad
\chi^2 = \chi^2_{\mathrm{CC}} + \chi^2_{\mathrm{SN}}.
\end{equation}
The H($z$) (CC+BAO) term is diagonal,
$\chi^2_{\mathrm{CC}}=\sum_i\big((H_{\mathrm{obs},i}-H(z_i;\theta))/\sigma_i\big)^2$.
The supernova term marginalizes the unknown absolute magnitude
$M_{\mathrm{abs}}$ analytically.

\begin{proposition}[Analytic $M_{\mathrm{abs}}$ marginalization, Goliath 2001]
With residuals $\Delta_i=m_{\mathrm{obs},i}-\mu_{\mathrm{th}}(z_i;\theta)$ and a
flat prior on the constant offset $M_{\mathrm{abs}}$, integrating it out gives
\begin{equation}
\chi^2_{\mathrm{eff}} = A - \frac{B^2}{C},\quad
A=\Delta^\top C_{\!*}^{-1}\Delta,\;
B=\Delta^\top C_{\!*}^{-1}\mathbf{1},\;
C=\mathbf{1}^\top C_{\!*}^{-1}\mathbf{1},
\end{equation}
with $C_{\!*}$ the (diagonal for Pantheon 2018, full for Pantheon+) covariance.
The result is invariant under $m_{\mathrm{obs}}\!\to\!m_{\mathrm{obs}}+\text{const}$.
\end{proposition}

The shift-invariance is the defining correctness property and is asserted in
\code{\seqsplit{test\_mabs\_marginalization\_shift\_invariance}}.

\begin{remark}[Combined datasets and $H_0$]
SNe with $M_{\mathrm{abs}}$ marginalized constrain only the \emph{shape}
$\Ez$, not $H_0$. When CC+BAO is added, $H_0$ is anchored \emph{directly} by
the H($z$) term: the combined fit does not marginalize $H_0$, only the SNe
nuisance offset. This is intentional --- CC+BAO is what pins $H_0$ --- and is
documented in \code{Posterior.chi2}.
\end{remark}

\subsection{Vectorized batch evaluation}

\code{log\_prob\_batch} evaluates the posterior for a batch of $B$ parameter
vectors in one pass, broadcasting $\Ezsq$ over the redshift grid. This is the
single hot path for the QVMC grid target and the vectorized QMCMC kernel.

\paragraph{Complexity.}
For a batch of $B$: the CC term is $\bigO(B N_{\mathrm{CC}})$; the SNe distance
integral is $\bigO(B N_z)$; the Pantheon+ quadratic form $\Delta^\top
C^{-1}\Delta$ is $\bigO(B N_{\mathrm{SN}}^2)$ (the $C^{-1}$ is precomputed
once). Time $\bigO\!\big(B(N_z+N_{\mathrm{SN}}^2)\big)$; space
$\bigO(B N_z + N_{\mathrm{SN}}^2)$.

\subsection{Diagnostics: $\tau$, ESS and split-$\Rhat$}

\begin{definition}[Integrated autocorrelation time]
For a stationary chain with normalized autocorrelation $\rho(t)$,
$\tau = 1 + 2\sum_{t\ge1}\rho(t)$. On a finite chain the sum must be truncated;
\code{autocorr\_time\_fft} uses \textbf{Sokal automatic windowing}: accumulate
$\tau(M)=1+2\sum_{t=1}^{M}\rho(t)$ and stop at the first $M$ with $M\ge c\,
\tau(M)$ ($c=5$). The autocovariance comes from one zero-padded FFT, so the
cost is $\bigO(N\log N)$.
\end{definition}

\begin{remark}[Why the old cutoff was wrong]
The previous rule ``first lag with $\rho<0.05$, else $N/4$'' returned the same
sentinel ($\arg\max=0$) for a chain that never decorrelates \emph{and} for one
already decorrelated at lag 1, conflating the worst and best chains. The Sokal
rule is self-consistent and is the community standard (emcee). Tests:
\code{test\_autocorr\_white\_noise} ($\tau\!\approx\!1$),
\code{test\_autocorr\_ar1} (AR(1) $\phi=0.9\Rightarrow\tau\!\approx\!19$), and
the regression \code{test\_autocorr\_excellent\_chain\_regression}.
\end{remark}

The effective sample size is $\mathrm{ESS}=MN/\tau_{\max}$ with $\tau_{\max}$
the worst per-chain-averaged $\tau$ across parameters (the slowest direction
governs independence). For weighted QVMC samples the Kish ESS
$(\sum w)^2/\sum w^2$ is used instead.

\begin{definition}[Rank-normalized split-$\Rhat$, Vehtari et al.\ 2021]
Split each chain in half (catching within-chain drift), rank-normalize the
pooled draws (robust to heavy tails), then form
$\Rhat=\sqrt{\big(\tfrac{N-1}{N}W+B/N\big)/W}$ from the between- ($B$) and
within-chain ($W$) variances; report the max over parameters. Convergence
threshold $\Rhat<1.01$.
\end{definition}

\begin{remark}[Why 1.01 matters here]
The project's central claim is ``quantum reproduces classical.'' Agreement
between two \emph{unconverged} chains proves nothing, so a lax threshold
(1.05) would make the comparison vacuous. The stricter rank-normalized
split-$\Rhat<1.01$ guards this; tests
\code{test\_split\_rhat\_converged\_vs\_drift} and
\code{\seqsplit{test\_split\_rhat\_catches\_within\_chain\_trend}} show it separates
converged from drifting chains and catches a trend the classical $\Rhat$
misses.
\end{remark}

% =============================================================================
\section{Quantum Computation: from one qubit to the circuits}
\label{sec:bridge}
% =============================================================================

This is the bridge. We start where the reader is --- one qubit on the Bloch
sphere --- and build up exactly the machinery the samplers use.

\subsection{One qubit, rotations, and the birth of probability}

A qubit state $\ket{\psi}=\cos\tfrac{\theta}{2}\ket{0}
+e^{i\varphi}\sin\tfrac{\theta}{2}\ket{1}$ is a point $(\theta,\varphi)$ on the
Bloch sphere. The rotation $R_Y(\lambda)$ takes $\ket{0}$ to
$\cos\tfrac{\lambda}{2}\ket{0}+\sin\tfrac{\lambda}{2}\ket{1}$, so measuring in
the computational basis gives
\begin{equation}
P(\ket{1}) = \sin^2(\lambda/2),\qquad P(\ket{0})=\cos^2(\lambda/2).
\label{eq:ry}
\end{equation}
Equation~\eqref{eq:ry} is the entire trick behind the \faithful{} acceptance
cell below: any target probability $A\in[0,1]$ is obtained with
$\lambda=2\arccos\sqrt{A}$.

\subsection{Amplitude encoding of the Metropolis rule (a \faithful{} cell)}
\label{ssec:accept}

\begin{proposition}[Quantum acceptance reproduces Metropolis exactly]
Let $\Delta=\log P(\theta')-\log P(\theta)$ and $A=\min(1,e^{\Delta})$ the
Metropolis acceptance. Preparing $R_Y\!\big(2\arccos\sqrt{A}\big)\ket{0}$ and
reading $|\!\braket{0|\psi}\!|^2$ returns $A$ exactly. Hence a chain using the
quantum acceptance is identical (given the same proposal and random stream) to
the classical one.
\end{proposition}

This is why the acceptance is labelled \faithful{}: the circuit computes the
\emph{same number} the CPU would. \code{hadamard\_accept\_log\_batch} evaluates
all $M$ chains' single-qubit circuits in one Aer job; the test
\code{test\_qpu\_helpers::test\_b1\_metropolis\_rule} checks monotonicity in
$\Delta$ (the regression guard for the historically inverted readout).

\begin{center}
\begin{quantikz}
\lstick{$\ket{0}$} & \gate{R_Y(2\arccos\sqrt{A})} & \meter{} 
\end{quantikz}
\end{center}

\subsection{Registers, superposition and the parameter grid}

Stacking $n$ qubits and applying $H^{\otimes n}$ creates a uniform
superposition over $2^n$ basis states. The QVMC/QGA encode a discretized
parameter grid in such a register: parameter $i$ occupies $\mathrm{nqpp}$
qubits, so $d$ parameters need $d\cdot\mathrm{nqpp}$ qubits and the grid has
$2^{d\cdot\mathrm{nqpp}}$ points. \code{\_build\_theta\_table} is the bijection
between a basis-state index and a $\theta$ on this grid.

\subsection{Entanglement and the hardware-efficient ansatz}

Two-qubit gates (CX, controlled-$R_Y$) create entanglement --- correlations no
product state can represent. The QVMC \emph{ansatz} interleaves single-qubit
$R_Y/R_Z$ rotations with a ring of CX gates; its tunable angles $\varphi$
define a variational family $\ket{\psi(\varphi)}$ whose measurement
distribution $|\!\braket{x|\psi(\varphi)}\!|^2$ is fitted to the posterior.

\subsection{The quantum proposal (an \algorithmic{} cell)}

The QMCMC proposal circuit (Sarracino et al.\ 2025) is
$H^{\otimes n}\!\to[R_Y R_Z\text{ + chained }CR_Y]^{\times L}\!\to H^{\otimes
n}\!\to R_Y$; the real parts of its output amplitudes define a displacement of
$\theta$. Because the displacement depends only on random angles (not on the
current $\theta$), a whole block is generated in one Aer job and consumed from
a queue (\code{QuantumProposalEngine}). The raw displacement is zero-mean but
has std $\approx0.35$; it is rescaled to unit std so it is a drop-in for the
classical $\mathcal N(0,1)$ proposal (keeping acceptance in a healthy band).
This is an \algorithmic{} cell: it is genuinely not the classical Gaussian, so
it changes the chain.

\paragraph{Complexity.}
A statevector refill of $b$ proposals on $n_q$ qubits is
$\bigO(b\cdot 2^{n_q})$ time, $\bigO(2^{n_q})$ space; amortized per step it is
one queue pop, $\bigO(1)$.

% =============================================================================
\section{MCMC and Variational Inference: two ways to know a posterior}
\label{sec:mcmc_vi}
% =============================================================================

Both samplers attack the same object --- the posterior $P(\theta\mid
\mathrm{data})$ --- but in fundamentally different ways. Understanding why they
differ is the key to reading the results, and to understanding why some quantum
substitutions change the answer and others cannot.

\subsection{MCMC: approximate by \emph{sampling}}

Markov Chain Monte Carlo does not try to write down $P$. It builds a random
walk through parameter space whose \emph{visited points}, after equilibration,
are distributed according to $P$. The Metropolis--Hastings recipe: from the
current $\theta_t$, draw a proposal $\theta'\sim q(\cdot\mid\theta_t)$, and
accept it with probability
\begin{equation}
A(\theta_t\!\to\!\theta') = \min\!\Big(1,\;
\frac{P(\theta')\,q(\theta_t\mid\theta')}{P(\theta_t)\,q(\theta'\mid\theta_t)}\Big).
\label{eq:mh}
\end{equation}
For a symmetric proposal ($q(\theta'\mid\theta_t)=q(\theta_t\mid\theta')$, e.g.
a Gaussian step) the ratio collapses to $A=\min(1,P(\theta')/P(\theta_t))=
\min(1,e^{\Delta})$ with $\Delta=\log P(\theta')-\log P(\theta_t)$. Accepted
and rejected steps both produce a sample (on rejection $\theta_{t+1}=\theta_t$).

\begin{proposition}[Why it works]
The Metropolis rule satisfies \emph{detailed balance},
$P(\theta)\,T(\theta\!\to\!\theta')=P(\theta')\,T(\theta'\!\to\!\theta)$, where
$T$ is the full transition kernel (proposal $\times$ acceptance). Detailed
balance makes $P$ the stationary distribution of the chain, so for large $t$
the samples $\{\theta_t\}$ are draws from $P$ --- regardless of the proposal
$q$, which affects only the \emph{efficiency} of mixing, not the
\emph{destination}.
\end{proposition}

The two sub-algorithms of MCMC are therefore:
\begin{itemize}[nosep]
\item \textbf{proposal} --- where to step ($q$); changes how fast you explore,
not where you end up;
\item \textbf{acceptance} --- whether to keep the step (Eq.~\ref{eq:mh});
enforces detailed balance and fixes the stationary distribution.
\end{itemize}
The output is a set of correlated samples; means, credible intervals and the
corner plot are computed by histogramming them. Diagnostics ($\tau$, ESS,
split-$\Rhat$, \S\ref{sec:stats}) ask ``has the walk equilibrated and how many
\emph{independent} samples do I effectively have?''

\subsection{Variational Inference: approximate by \emph{optimizing}}

VI takes the opposite stance. Instead of sampling, it picks a family of simple
distributions $Q_\varphi(\theta)$ and \emph{tunes} $\varphi$ until $Q_\varphi$
is as close as possible to $P$, measured by the Kullback--Leibler divergence
\begin{equation}
\KL(Q_\varphi\,\|\,P) = \int Q_\varphi(\theta)\,
\log\frac{Q_\varphi(\theta)}{P(\theta\mid\mathrm{data})}\,d\theta
= \mathbb{E}_{Q_\varphi}\!\big[\log Q_\varphi - \log P\big].
\label{eq:kl}
\end{equation}
Inference becomes \emph{optimization}: $\varphi^\star=\arg\min_\varphi
\KL(Q_\varphi\|P)$, solved by gradient descent (classical VI here uses a grid
representation of $Q_\varphi$ and COBYLA/SGD). Once trained, $Q_{\varphi^\star}$
\emph{is} the approximate posterior --- you have a closed object you can sample
or read moments from directly, with no autocorrelation and no burn-in.

In this project $Q_\varphi$ lives on the discrete $2^{\mathrm{nqpp}}$ grid per
parameter (\S\ref{sec:bridge}), so the two sub-algorithms of VI are:
\begin{itemize}[nosep]
\item \textbf{training} --- minimize Eq.~\ref{eq:kl} over $\varphi$; this is
where the approximation is actually fitted;
\item \textbf{sampling} --- once trained, draw from $Q_{\varphi^\star}$ to
report means and a corner plot;
\item \textbf{normalization} --- ensure $Q_\varphi$ sums to one over the
(masked) grid support, so the KL and the reported distribution are proper.
\end{itemize}

\subsection{Why they are fundamentally different}

\begin{table}[h]
\centering\footnotesize
\begin{tabular}{@{}lll@{}}
\toprule
 & \textbf{MCMC} & \textbf{Variational Inference}\\
\midrule
Strategy & sample the posterior & optimize a proxy to the posterior\\
Output & correlated draws & a fitted closed-form $Q_{\varphi^\star}$\\
Exactness & asymptotically exact (given mixing) & only as good as the family $Q_\varphi$\\
Failure mode & slow mixing / non-convergence & too-rigid family, biased $Q$\\
Cost driver & \#steps $\times$ likelihood & \#iters $\times$ grid $2^{n\cdot\mathrm{nqpp}}$\\
Diagnostic & $\Rhat$, ESS, $\tau$ & KL curve vs iteration\\
\bottomrule
\end{tabular}
\caption{MCMC is unbiased but can be slow to converge; VI is fast and gives a
closed object but is only as good as the chosen family $Q_\varphi$ can
represent. They are complementary, not interchangeable --- which is why the
project runs both and compares them against the \emph{same} classical
baselines.}
\end{table}

A blunt way to put it: MCMC is asymptotically exact but you never quite know
when you have arrived; VI always ``arrives'' (the optimizer converges) but to a
point that is only as correct as your family $Q_\varphi$ allows. The KL floor
you see in the QVMC training curve (it bottoms out at some $\KL>0$) is exactly
this family-limitation: even a perfectly trained $Q_\varphi$ on a coarse grid
cannot match $P$ to zero KL.

\subsection{The quantum counterparts, sub-algorithm by sub-algorithm}

Each classical sub-algorithm has a quantum re-encoding. The crucial point ---
and the heart of the ablation framework (\S\ref{sec:ablation}) --- is that some
re-encodings compute the \emph{identical} mathematical object (\faithful), while
others are genuinely different algorithms (\algorithmic).

\begin{center}
\footnotesize
\begin{longtable}{@{}p{2.1cm}p{4.3cm}p{4.6cm}c@{}}
\toprule
\textbf{Sub-alg.} & \textbf{Classical} & \textbf{Quantum} & \textbf{Kind}\\
\midrule
\endhead
MCMC proposal & Gaussian step $\theta'=\theta+\mathcal N(0,\sigma)$ &
Statevector circuit (Sarracino): random-angle circuit whose amplitudes define
the displacement, rescaled to unit std & \algorithmic\\[4pt]
MCMC acceptance & $A=\min(1,e^\Delta)$ on the CPU &
$R_Y(2\arccos\!\sqrt{A})\ket0$, read $P(\ket0)=A$ --- the \emph{same} number &
\faithful\\[4pt]
VI training & gradient descent on $\KL$ (COBYLA/SGD) &
parameter-shift gradients of a variational circuit (or SPSA on hardware) ---
a different optimizer reaching a different optimum & \algorithmic\\[4pt]
VI sampling & draw indices from the trained grid pmf &
measure shots of the trained state $|\psi(\varphi^\star)|^2$ --- same
distribution, finite-shot noise & \faithful\\[4pt]
VI normalization & divide by the sum over support &
exact renormalization over the masked support --- identical & \faithful\\
\bottomrule
\end{longtable}
\end{center}

\subsection{Why QMCMC 50\% and 100\% give the same corner plot}
\label{ssec:why5050}

This is the question the corner plots provoke, and it has a clean answer.

Recall the QMCMC ladder switches its two sub-algorithms to quantum one at a
time:
\[
\underbrace{0\%}_{\text{classical}}
\;\to\;
\underbrace{50\%}_{\text{+ quantum \emph{proposal}}}
\;\to\;
\underbrace{100\%}_{\text{+ quantum \emph{acceptance}}}.
\]
The step from 50\% to 100\% switches \textbf{only the acceptance} from the CPU
formula $\min(1,e^\Delta)$ to its quantum encoding $R_Y(2\arccos\sqrt A)\ket0$.
But by the proposition of \S\ref{ssec:accept}, that circuit returns
\emph{exactly} $A=\min(1,e^\Delta)$ --- the same number, to machine precision.
The acceptance is a \faithful{} cell.

Now recall \emph{why} the acceptance cannot change the answer even in
principle: by detailed balance, the acceptance rule \emph{defines the
stationary distribution}, and both 50\% and 100\% use the \emph{same}
acceptance rule (Metropolis), computed two different ways that agree bit-for-
bit. With the same proposal, the same rule, and the same random-number stream,
the two chains take \emph{identical} steps. Their samples are identical, so
their corner plots, means and credible intervals coincide exactly.

Contrast the $0\%\!\to\!50\%$ step, which switches the \emph{proposal}
(\algorithmic): the quantum proposal is a genuinely different $q$, so the chain
explores differently and the corner plot shifts slightly (same destination $P$,
different path and mixing). And contrast the QVMC ladder, where the
$33\%\!\to\!67\%$ step switches \emph{training} (\algorithmic) and the KL drops
sharply --- there the quantum sub-algorithm really is a different optimizer.

\begin{remark}[The take-away]
A coincident corner plot between two adjacent rungs is not a bug and not a
boring null result: it is a \emph{measurement}. It confirms that the quantum
re-encoding of that sub-algorithm is faithful --- it reproduces the classical
computation exactly. That confirmation is the whole point of the project's
``fidelity, not advantage'' thesis. The rungs that \emph{should} coincide do,
and the rungs that \emph{should} differ do.
\end{remark}


% =============================================================================
\section{The samplers}
% =============================================================================

\subsection{QMCMC}
\label{sec:qmcmc}

\code{QMCMCModular} is a hand-written Metropolis--Hastings (not emcee), so that
each component can be swapped classical$\leftrightarrow$quantum and the
baseline and quantum run share the exact transition structure, step scale and
RNG stream. The kernel is vectorized: one \code{log\_prob\_batch} call scores
all $M$ chains per step.

\paragraph{Complexity.}
Per step the cost is dominated by the likelihood, $\bigO(M\,c_{\mathrm{post}})$
with $c_{\mathrm{post}}=\bigO(N_z+N_{\mathrm{SN}}^2)$; the quantum proposal
adds $\bigO(M)$ queue pops and the quantum acceptance one $\bigO(M)$ Aer job.
Space $\bigO(MNd)$ for the stored chains.

\subsection{QVMC and the adaptive grid}
\label{sec:qvmc}

\code{QVMCModular} represents the posterior on the discrete grid as
$|\psi(\varphi)|^2$ and minimizes $\KL(Q_\varphi\,\|\,P)$. A naive grid over
the full \code{sample\_box} collapses the (narrow) cosmological posterior onto
1--3 cells; \code{estimate\_grid\_window} (now in \code{cosmo\_core}, shared
with the QPU path) runs a short classical pre-fit and centres a zoomed
$\pm k\sigma$ window on the posterior, using the \emph{median} (robust to
asymmetric posteriors).

Training is the \algorithmic{} cell: quantum parameter-shift gradients (or
SPSA on hardware) reach a different KL minimum than classical COBYLA. Sampling
and normalization are \faithful{} cells (shots from the trained state; exact
renormalization).

\paragraph{Complexity.}
Building the target evaluates the posterior on all $2^{d\cdot\mathrm{nqpp}}$
grid points: $\bigO(2^{d\cdot\mathrm{nqpp}}\,c_{\mathrm{post}})$ time and
$\bigO(2^{d\cdot\mathrm{nqpp}})$ space --- the \textbf{exponential} cost in
$d\cdot\mathrm{nqpp}$ that the \code{--max-qubits} guard protects.
Parameter-shift adds $2n_\varphi$ circuit evaluations per iteration, batched
into one Aer job.

\subsection{Genetic optimizers (CGA / QGA)}
\label{sec:genetic}

The genetic fitness is the \emph{same} log-posterior, so maximizing fitness
$\equiv$ finding the MAP. CGA is pure NumPy; QGA substitutes quantum operators
(init / mutation / crossover), all \algorithmic{}. 

\begin{remark}[The transpile-once fix, A1]
Each QGA operator circuit is \emph{parametrized} and transpiled once; an
individual's gene bits enter only through bound rotation angles (a bit $b$ with
flip probability $p$ maps to $R_Y(2\arcsin\sqrt{b(1{-}p)+(1{-}b)p})$). The hot
loop then only \code{assign\_parameters} on the cached transpiled template,
instead of building and transpiling one fresh circuit per individual per
generation ($P\!\cdot\!d$ transpilations/generation in the old code).
\end{remark}

\paragraph{Complexity.}
$\bigO(G\cdot P\cdot d)$ classical work per run; the QGA adds one batched Aer
job per operator per generation (no per-individual transpilation after A1).

% =============================================================================
\section{The classical--quantum ablation framework}
\label{sec:ablation}
% =============================================================================

\begin{definition}[Ablation index, uniform weighting]
The \emph{quantumness} of a configuration is the fraction of substitution
points switched to quantum, every point weighted equally:
\begin{equation}
\text{index} = 100\cdot\frac{\#\{\text{quantum components}\}}{\#\{\text{components}\}}.
\end{equation}
It is an \textbf{ablation index}, not a quantum-resource metric: it counts how
many algorithmic steps are quantum, nothing more.
\end{definition}

Each substitution point is \faithful{} (a quantum re-encoding of the exact
classical rule; a null cell whose job is to verify the substitution introduces
no bias) or \algorithmic{} (a genuinely different quantum algorithm; a
treatment cell). The per-method ladders (uniform spacing) are
\[
\text{QMCMC: } 0\to50\to100\,\%,\qquad
\text{QVMC: } 0\to33\to67\to100\,\%,\qquad
\text{QGA: } 0\to33\to67\to100\,\%.
\]

\begin{table}[h]
\centering\footnotesize
\begin{tabular}{@{}llll@{}}
\toprule
Sampler & Substitution point & Kind & Expected outcome of flipping it\\
\midrule
QMCMC & proposal      & \algorithmic & changes (genuine quantum proposal)\\
QMCMC & acceptance    & \faithful    & identical (reproduces Metropolis)\\
QVMC  & sampling      & \faithful    & $\approx$identical (shot noise only)\\
QVMC  & training      & \algorithmic & changes strongly (lower KL)\\
QVMC  & normalization & \faithful    & $\approx$identical (faithful renorm.)\\
\bottomrule
\end{tabular}
\caption{The ablation map. Faithful cells are falsifiable correctness tests;
algorithmic cells are where the science happens.}
\end{table}

\begin{remark}[Why this resolves the ``ornamental circuit'' objection]
A faithful cell runs a quantum circuit that, by construction, returns the same
number as the CPU. Far from being ornamental, it is the \emph{control} of the
experiment: its coincidence with the classical baseline is the evidence that
the quantum re-encoding is faithful. The framework thus turns the project's
narrative (``some rungs coincide by design, others differ'') into first-class,
testable metadata. The historical hand-picked weights are retained only as
\code{legacy\_weighted\_index} for CSV continuity.
\end{remark}

% =============================================================================
\section{Complexity summary}
\label{sec:complexity}
% =============================================================================

\begin{center}
\footnotesize
\begin{longtable}{@{}lll@{}}
\toprule
\textbf{Component} & \textbf{Time} & \textbf{Space}\\
\midrule
\endhead
$\Ezsq$, $H(z)$              & $\bigO(N_z)$ & $\bigO(N_z)$\\
$\mu_{\mathrm{th}}$ (cumulative trapezoid) & $\bigO(N_z)$ & $\bigO(N_z)$\\
\code{log\_prob\_batch}      & $\bigO(B(N_z+N_{\mathrm{SN}}^2))$ & $\bigO(BN_z+N_{\mathrm{SN}}^2)$\\
$\tau$ (FFT, Sokal)          & $\bigO(N\log N)$ & $\bigO(N)$\\
split-$\Rhat$                & $\bigO(MN)$ & $\bigO(MN)$\\
QMCMC (per step)             & $\bigO(M\,c_{\mathrm{post}})$ & $\bigO(MNd)$\\
quantum proposal (refill)    & $\bigO(b\,2^{n_q})$ & $\bigO(2^{n_q})$\\
quantum acceptance (batch)   & $\bigO(M)$ & $\bigO(M)$\\
QVMC target build            & $\bigO(2^{d\cdot\mathrm{nqpp}}c_{\mathrm{post}})$ & $\bigO(2^{d\cdot\mathrm{nqpp}})$\\
QVMC parameter-shift (iter)  & $\bigO(n_\varphi\,2^{d\cdot\mathrm{nqpp}})$ & $\bigO(2^{d\cdot\mathrm{nqpp}})$\\
CGA / QGA (per run)          & $\bigO(GPd)$ & $\bigO(Pd)$\\
\bottomrule
\end{longtable}
\end{center}

The single dominant scaling to internalize is the QVMC/QGA grid:
$2^{d\cdot\mathrm{nqpp}}$ in both time and space. This is the exponential wall
that \code{--max-qubits} guards and the reason a QPU --- which never
materializes the statevector --- is the only way past it for large grids,
though at the cost of shot noise and queue-dominated wall time rather than RAM.

\section*{Acknowledgements}
This document accompanies the \emph{Quantum Algorithms for Cosmology}
codebase. The scientific framing throughout is fidelity, not advantage.

\end{document}
